
\SGARefTarget{sga3:xvi:section:5}%
\section*{5. Application to Reductive and Semisimple Groups}
\addcontentsline{toc}{section}{5. Application to Reductive and
Semisimple Groups}

\noindent\SGARefTarget{sga3:xvi:definition:5:1}%
\textbf{Definition 5.1.} --- \textit{An \(S\)-group prescheme
\(G\) is said to be \emph{reductive} (respectively
\emph{semisimple}) if \(G\) is smooth and affine over \(S\), with
reductive (respectively semisimple) fibers.}

\medskip
\noindent\SGARefTarget{sga3:xvi:5:1:1}%
\textbf{5.1.1.}
Reductive groups will be studied systematically beginning with
Exposé XIX. In this section we shall need the following
properties, which will be proved in Exposé XIX (without using the
developments of the present Exposé).

\begin{enumerate}
\renewcommand{\labelenumi}{\textit{\alph{enumi})}}
\item \textit{Let \(S\) be a prescheme, let \(G\) be an
\(S\)-group prescheme smooth and affine over \(S\), with connected
fibers, and let \(s\) be a point of \(S\) such that \(G_s\) is
reductive (respectively semisimple). Then there exists a
neighborhood \(U\) of \(s\) such that \(G|_U\) is reductive
(respectively semisimple).}

\item \textit{If \(G\) is reductive, then locally for the étale
topology \(G\) possesses maximal tori; and if \(T\) is a maximal
torus of \(G\), the Weyl group}
\[
  W=\Norm_G(T)/T
\]
\textit{is finite over \(S\).}

\item \textit{A reductive group has unipotent rank zero.}
\end{enumerate}

Admitting these facts, we propose to strengthen assertion a)
above.

\medskip
\noindent\SGARefTarget{sga3:xvi:theorem:5:2}%
\textbf{Theorem 5.2.} --- \textit{Let \(S\) be a prescheme, let
\(G\) be an \(S\)-group prescheme smooth and of finite
presentation over \(S\), with connected fibers, and let \(s\) be
a point of \(S\). Then:}

\begin{enumerate}
\renewcommand{\labelenumi}{\textit{(\roman{enumi})}}
\item \textit{If the fibers of \(G\) are affine and \(G_s\) is
reductive, there exists an open neighborhood \(U\) of \(s\) such
that \(G|_U\) is reductive.}

\item \textit{If \(G_s\) is semisimple, there exists an open
neighborhood \(U\) of \(s\) such that \(G|_U\) is semisimple.}
\end{enumerate}

Using \SGARefLink{sga3:xv:proposition:6-2}{Exposé XV, 6.2} i), and Exposé VI\textsubscript{B}, \S10, we
reduce to the case where \(S\) is noetherian.

In case ii), consider the center \(Z\) of \(G\), which is
representable (\SGARefLink{sga3:xi:corollary:6:11}{Exposé XI, 6.11} a)). Since \(G_s\) is semisimple,
it is well known that \(Z_s\) is finite. Consequently (Exposé
VI\textsubscript{B}, \S4), there exists a neighborhood \(U\) of
\(s\) such that \(Z\) is quasi-finite over \(U\). For every point
\(t\) of \(U\), \(G_t\) is then affine (\SGARefLink{sga3:xii:6.1}{Exposé XII, 6.1}).
Shrinking \(S\) if necessary, we may therefore suppose that the
fibers of \(G\) are affine, in case ii) as well as in case i).

In view of a), to prove \SGARefLink{sga3:xvi:theorem:5:2}{5.2} it suffices to prove that \(G\) is
affine over \(S\). Since \(G\) has affine fibers, we know that the
functor of subtori of \(G\) is representable by a smooth
\(S\)-prescheme (\SGARefLink{sga3:xv:corollary:8:11}{Exposé XV, 8.11} and \SGARefLink{sga3:xv:theorem:8:8}{8.9}). After making an étale
extension covering \(s\), which is permissible, we may therefore
suppose that there is a subtorus \(T\) of \(G\) such that \(T_s\)
is a maximal torus of \(G_s\). But then
\[
  C=\Centr_G(T)
\]
is a smooth subgroup of \(G\), with connected fibers, containing
\(T\), such that \(C_s=T_s\) (because \(C_s\) is a Cartan
subgroup of \(G_s\), and \(G_s\) has unipotent rank zero by c)).
It follows that \(C=T\), so \(T\) is a Cartan subgroup and a
maximal torus of \(G\); in particular, the unipotent rank of \(G\)
is zero, and \SGARefLink{sga3:xvi:theorem:4:1}{4.1} applies.

By \SGARefLink{sga3:xvi:theorem:4:1}{4.1} d), it suffices to show that the Weyl group \(W\) relative
to \(T\) is affine over a neighborhood \(U\) of \(s\). In fact,
we shall show that \(W\) is even finite over a neighborhood of
\(s\). Since \(W\) is quasi-finite over \(S\) (\SGARefLink{sga3:xv:proposition:7-1}{Exposé XV,
7.1} iv)), saying that \(W\) is finite over a neighborhood \(U\)
of \(s\) is equivalent to saying that the morphism
\[
  W\longrightarrow S
\]
is proper at \(s\) (EGA IV, 15.7, and EGA III, 4.4.2). To
establish this, the valuative criterion of local properness
(EGA IV, 15.7) reduces us to the case where \(S\) is the spectrum
of a discrete valuation ring, with closed point \(s\). But then,
by the last assertion of \SGARefLink{sga3:xvi:theorem:4:1}{4.1} d), \(G\) is affine over \(S\), and
we may apply property a) recalled above to conclude that \(G\) is
reductive (respectively semisimple). Applying property b), we
conclude that \(W\) is indeed finite over \(S\), which completes
the proof of \SGARefLink{sga3:xvi:theorem:5:2}{5.2}.
